\section*{Requerimientos del Producto (PRD) -- Collar DR-VF-01}

\begin{table}[H]
\centering
\small
\begin{tabularx}{\textwidth}{|c X c|}
\hline
\rowcolor{categorybg}
\textbf{ID} & \textbf{Descripción del requerimiento} & \textbf{Tipo} \\ \hline
F-01 & El sistema debe determinar la ubicación geográfica mediante GPS u-blox MAX-M10S con actualización $>$ 1 Hz. & Funcional \\ \hline
F-02 & El microcontrolador ESP32-S3 debe ejecutar algoritmos Haversine y Point-in-Polygon en tiempo real. & Funcional \\ \hline
F-03 & El dispositivo debe contar con disuasión en dos etapas: alerta acústica y pulso disuasivo. & Funcional \\ \hline
F-04 & El sistema debe contar con aislamiento galvánico mediante optoacoplador (VO617A / PC817). & Funcional \\ \hline
F-05 & El collar debe contar con conectividad LoRa mediante el módulo DX-LR30-900M22S (SX1262, 915 MHz). (Módem celular proyectado en Fase 2; modelo \textbf{POR DECIDIR}: BG95-M3 o A7670SA). & Funcional \\ \hline
F-06 & Configuración de geocerca mediante servidor web interno accesible por Wi-Fi/BLE (LittleFS). & Funcional \\ \hline
F-07 & Carga solar conmutada MPPT mediante controlador CN3791 (6V) alimentado solo desde el panel solar, y carga por USB-C mediante el cargador lineal TP4056 (500 mA); el USB no alimenta el sistema. & Funcional \\ \hline
F-08 & Gestión de energía con lectura de batería relativa (0-100\%) mediante medidor I2C MAX17048. & Funcional \\ \hline
F-09 & El sistema debe regular la lógica de 3,3 V con un LDO TLV75733PDBVR (1 A, corriente quiescente de 25 $\mu$A) alimentado solo desde la batería. & Funcional \\ \hline
F-10 & El sistema debe proteger la batería LiPo con un circuito integrado de protección (\textbf{POR DECIDIR}; la celda elegida puede traer protección propia). & Funcional \\ \hline
F-11 & El puerto USB-C debe contar con protección contra descargas electrostáticas en D+, D- y VBUS mediante USBLC6-2SC6. & Funcional \\ \hline
F-12 & El microcontrolador debe permanecer en Deep Sleep y despertar por la interrupción INT1 de la IMU (GPIO 4) ante movimiento significativo (IMU \textbf{POR DECIDIR}: LSM6DSOX o LSM6DS3TR-C). & Funcional \\ \hline
NF-01 & Autonomía mínima de 7 a 10 días continuos en ausencia total de radiación solar (estimación teórica: unos 190 días con 1200 mAh, calculada en \texttt{hardware/POWER\_BUDGET.md}). & No funcional \\ \hline
NF-02 & Aislamiento galvánico con tensión de ruptura $\ge 3000~V_{RMS}$ entre control y etapa disuasiva. & No funcional \\ \hline
NF-03 & Encapsulado mecánico para prototipo con protección contra salpicaduras (IP67 en Fase 2). & No funcional \\ \hline
NF-04 & Costo total de componentes de la placa (BOM Target) $<$ \$150.000 COP (~35-40 USD). & No funcional \\ \hline
NF-05 & Rango térmico de operación expuesto en campo de $-10^\circ C$ a $50^\circ C$. & No funcional \\ \hline
NF-06 & Peso total del dispositivo y chasis de pruebas $\le$ 400 g. & No funcional \\ \hline
NF-07 & Coordenadas del perímetro guardadas en memoria no volátil (Flash NVS). & No funcional \\ \hline
NF-08 & PCB diseñada como Carrier Board a 2 capas en Altium Designer / EasyEDA. & No funcional \\ \hline
\end{tabularx}
\end{table}